\begingroup\let\clearpageforsection\relax\exercisesheader{}\endgroup

% 29

\eoce{\qt{Kondisi pencahayaan dan kinerja ujian\label{light_exam_performance}} Sebuah studi 
dirancang untuk menguji pengaruh kondisi pencahayaan terhadap kinerja ujian mahasiswa. 
Peneliti menduga bahwa kondisi pencahayaan dapat memberikan pengaruh yang berbeda pada 
laki-laki dan perempuan, sehingga keduanya hendak diwakili sama banyak dalam setiap 
perlakuan. Perlakuannya adalah pencahayaan fluoresen dari atas, pencahayaan kuning dari 
atas, dan tanpa pencahayaan dari atas (hanya menggunakan lampu meja). 
\begin{parts}
\item Apa variabel responsnya?
\item Apa variabel penjelasnya? Apa saja levelnya?
\item Apa variabel pemblokannya? Apa saja levelnya?
\end{parts}
}{}

% 30

\eoce{\qt{Suplemen vitamin\label{vitamin_supplement}}
Untuk menilai efektivitas konsumsi vitamin C dosis tinggi dalam mengurangi durasi selesma, 
para peneliti merekrut 400 sukarelawan sehat dari kalangan staf dan mahasiswa di sebuah 
universitas. Seperempat peserta dialokasikan untuk menerima plasebo, sedangkan peserta 
lainnya dialokasikan secara merata ke tiga perlakuan: vitamin C dosis 1~g, vitamin C dosis 
3~g, atau vitamin C dosis 3~g ditambah zat aditif. Pil tersebut harus mulai diminum ketika 
selesma muncul dan dilanjutkan selama dua hari berikutnya. Semua tablet memiliki tampilan 
dan kemasan yang identik. Perawat yang menyerahkan pil yang diresepkan mengetahui 
perlakuan yang diterima setiap peserta, tetapi 
peneliti yang menilai peserta ketika mereka sakit tidak mengetahuinya. Tidak ditemukan 
perbedaan signifikan di antara keempat kelompok pada ukuran durasi maupun tingkat 
keparahan selesma, dan kelompok plasebo mengalami durasi gejala paling 
singkat.\footfullcite{Audera:2001}
\begin{parts}
\item Apakah studi ini merupakan eksperimen atau studi observasional? Mengapa?
\item Apa variabel penjelas dan variabel respons dalam studi ini?
\item Apakah identitas perlakuan disamarkan dari peserta?
\item Apakah studi ini tersamar ganda?
\item Peserta pada akhirnya dapat memilih apakah akan meminum pil yang diresepkan. Kita 
dapat menduga bahwa tidak semuanya akan mematuhi penugasan dan meminum pilnya. Apakah 
ketidakpatuhan ini dengan sendirinya menciptakan variabel perancu? Jelaskan mengapa alokasi 
awal tetap acak, serta bagaimana ketidakpatuhan dapat mempersulit estimasi efek perlakuan 
yang benar-benar diterima dan membiaskan analisis yang mengelompokkan peserta menurut pil 
yang benar-benar mereka minum.
\end{parts}
}{}

% 31

\eoce{\qt{Kondisi pencahayaan, kebisingan, dan kinerja ujian\label{light_noise_exam_performance}} 
Sebuah studi dirancang untuk menguji pengaruh kondisi pencahayaan dan tingkat kebisingan 
terhadap kinerja ujian mahasiswa. Peneliti menduga bahwa kedua faktor tersebut dapat 
memberikan pengaruh yang berbeda pada laki-laki dan perempuan, sehingga keduanya hendak 
diwakili sama banyak dalam setiap perlakuan. Perlakuan pencahayaan yang dipertimbangkan 
adalah pencahayaan fluoresen dari atas, pencahayaan kuning dari atas, dan tanpa pencahayaan 
dari atas (hanya menggunakan lampu meja). Perlakuan kebisingannya adalah tanpa kebisingan, 
kebisingan konstruksi, dan suara percakapan manusia.
\begin{parts}
\item Apa jenis studi ini?
\item Berapa faktor yang dipertimbangkan dalam studi ini? Sebutkan faktor-faktor tersebut 
dan uraikan levelnya.
\item Apa peran variabel jenis kelamin dalam studi ini?
\end{parts}
}{}

% 32

\eoce{\qt{Musik dan pembelajaran\label{music_learning}} Anda hendak melakukan eksperimen di 
kelas untuk mengetahui apakah mahasiswa belajar lebih baik ketika belajar tanpa musik, 
dengan musik tanpa lirik (instrumental), atau dengan musik berlirik. Uraikan secara 
ringkas sebuah rancangan untuk studi ini.
}{}

% 33

\eoce{\qt{Preferensi minuman bersoda\label{soda_preference}} Anda hendak melakukan eksperimen 
di kelas untuk mengetahui apakah teman-teman sekelas Anda lebih menyukai rasa Coke biasa 
atau Diet Coke. Uraikan secara ringkas sebuah rancangan untuk studi ini.
}{}

% 34

\eoce{\qt{Olahraga dan kesehatan mental\label{exercise_mental_health}} Seorang peneliti tertarik 
pada pengaruh olahraga terhadap kesehatan mental dan mengusulkan studi berikut: Gunakan 
pengambilan sampel acak berstrata untuk memastikan bahwa proporsi kelompok usia 18--30, 
31--40, dan 41--55 tahun dalam sampel mewakili populasi. Selanjutnya, alokasikan secara 
acak separuh subjek dalam setiap kelompok usia untuk berolahraga dua kali seminggu, dan 
instruksikan subjek lainnya agar tidak berolahraga. Lakukan pemeriksaan kesehatan mental 
pada awal dan akhir studi, lalu bandingkan hasilnya.
\begin{parts}
\item Apa jenis studi ini? 
\item Apa kelompok perlakuan dan kelompok kontrol dalam studi ini?
\item Apakah studi ini menerapkan pemblokan? Jika ya, apa variabel pemblokannya?
\item Apakah studi ini menerapkan penyamaran?
\item Jelaskan apakah hasil studi dapat digunakan untuk menetapkan hubungan sebab-akibat 
antara olahraga dan kesehatan mental, serta apakah kesimpulannya dapat digeneralisasikan 
ke populasi secara luas.
\item Andaikan Anda ditugaskan untuk menentukan apakah studi yang diusulkan ini layak 
menerima pendanaan. Apakah Anda memiliki keberatan terhadap usulan studi tersebut?
\end{parts}
}{}
